\documentclass[12pt]{article}
\usepackage[utf8]{inputenc}
\usepackage[T1]{fontenc}
\usepackage{hyperref}
\title{Tasmanian Language Resources for K-12 Curriculum Development}
\author{Compiled by Claude Sonnet 5.5 (claude-sonnet-5-5) for Andr\'as Kornai}
\date{}
\begin{document}
\maketitle

\section{Introduction}
The traditional Tasmanian languages (ISO 639-3 xtz, Glottocode tasm1247) are no longer spoken. The identification given for this entry lists the family as ``Bookkeeping'', which is a Glottolog placeholder and not a genealogical family. The search results mostly concerned palawa kani, the revived composite language run by the Tasmanian Aboriginal Centre. Wikipedia lists palawa kani under Glottolog pala1356 and does not give it the code xtz. I did not verify how xtz relates to palawa kani, and the two should not be assumed to be the same thing. This survey rests on one web search.

\section{Primary Educational Resources}
No school textbooks or primers openly available were found. The sources report that dictionaries and other learning resources for palawa kani are restricted to the Aboriginal community. These resources are not publicly accessible.

\section{Linguistic and Reference Materials}
\begin{itemize}
\item Wikipedia, palawa kani (overview, also mentions Plomley's 1976 word lists and Crowley and Dixon 1981 as main historical sources). \url{https://en.wikipedia.org/wiki/Palawa_kani}
\end{itemize}

\section{Religious and Community Resources}
\begin{itemize}
\item Wukalina Walk, language guide. \url{https://www.wukalinawalk.com.au/language-guide}
\item Whose Knowledge podcast, Theresa Sainty and the path to revive palawa kani. \url{https://podcast.whoseknowledge.org/posts/ep12/}
\end{itemize}

\section{Digital and Audio Resources}
\begin{itemize}
\item Ionata, ``Advancing Language Research and Revival'' (article reporting a place names app with over 300 names; it is promotional). \url{https://ionata.com.au/2025/12/advancing-language-research-and-revival/}
\end{itemize}

\section{Curriculum Development Context}
TasTAFE links to the Tasmanian Aboriginal Centre's palawa kani program and its policy and protocol for using the language: \url{https://www.tastafe.tas.edu.au/students/student-services/taypani}. Any curriculum work would have to go through the Tasmanian Aboriginal Centre and follow its protocols. Little open material was found.
\end{document}